\documentclass[../DoAn.tex]{subfiles}
\begin{document}
\label{ch:co_so}
% =============================================================
% CHƯƠNG 2 — NỀN TẢNG LÝ THUYẾT          Mục tiêu độ dài: ~10 trang
% =============================================================

Chương này hệ thống hoá các kiến thức nền tảng phục vụ cho phương pháp đề xuất ở Chương~\ref{ch:de_xuat}. Phần đầu trình bày mô hình hoá bài toán nhúng mạng ảo và phát biểu dưới dạng quy hoạch nguyên tuyến tính. Tiếp theo là khảo sát các hướng tiếp cận hiện có cùng phân tích ưu nhược. Phần sau trình bày các công cụ học tăng cường sâu và mạng nơ-ron đồ thị, là hai khối nền tảng được sử dụng trong đồ án.

\section{Mô hình hoá bài toán nhúng mạng ảo}
\label{sec:vne_model}

\subsection{Mạng vật lý và mạng ảo}

\textbf{Mạng vật lý} (substrate network) được biểu diễn bằng đồ thị vô hướng có trọng số
\[
    G^s = (N^s, E^s),
\]
trong đó $N^s$ là tập các nút vật lý và $E^s$ là tập các liên kết. Mỗi nút $n^s \in N^s$ được đặc trưng bởi dung lượng CPU $C(n^s)$, đơn giá CPU $P(n^s)$ và độ trễ xử lý $D(n^s)$. Mỗi liên kết $e^s = (u^s, v^s) \in E^s$ được đặc trưng bởi băng thông $B(e^s)$, đơn giá băng thông $P(e^s)$ và độ trễ truyền dẫn $D(e^s)$. Tại thời điểm bất kỳ, tài nguyên còn lại của mỗi nút và liên kết là $C^{\text{avail}}(n^s) \leq C(n^s)$ và $B^{\text{avail}}(e^s) \leq B(e^s)$.

Một \textbf{yêu cầu mạng ảo} (Virtual Network Request - VNR) bao gồm đồ thị mạng ảo và các thông tin thời gian:
\[
    \text{VNR} = (G^v, t_a, t_l), \quad G^v = (N^v, E^v),
\]
với $t_a$ là thời điểm đến và $t_l$ là thời gian sống. Mỗi nút ảo $n^v \in N^v$ có yêu cầu CPU $c(n^v)$, và mỗi liên kết ảo $e^v = (u^v, v^v) \in E^v$ có yêu cầu băng thông $b(e^v)$. Trong cấu hình đa miền, mỗi nút ảo có thể có thêm tập miền cho phép $\mathcal{D}(n^v) \subseteq \mathcal{D}$.

\subsection{Hai pha của VNE: ánh xạ nút và ánh xạ liên kết}

Một lời giải nhúng cho VNR là một cặp ánh xạ $(f_N, f_E)$ trong đó:
\begin{itemize}
\item \textbf{Ánh xạ nút} $f_N : N^v \to N^s$ thoả mãn (i) ràng buộc tài nguyên $c(n^v) \leq C^{\text{avail}}(f_N(n^v))$ với mọi $n^v$, (ii) ràng buộc đơn ánh: $f_N(u^v) \neq f_N(v^v)$ với $u^v \neq v^v$, và (iii) ràng buộc miền $f_N(n^v) \in \mathcal{D}(n^v)$ nếu có.
\item \textbf{Ánh xạ liên kết} $f_E : E^v \to \mathcal{P}(E^s)$ ánh xạ mỗi liên kết ảo tới một tập đường đi vật lý nối $f_N(u^v)$ và $f_N(v^v)$, với tổng băng thông trên các đường đi bằng $b(e^v)$ và không vượt quá băng thông còn lại của bất kỳ liên kết vật lý nào.
\end{itemize}

Do tính chất NP-Hard của bài toán, các phương pháp thực tế thường giải hai pha một cách phối hợp nhưng không đồng thời tối ưu: pha ánh xạ nút sử dụng một chính sách xếp hạng/lựa chọn, sau đó pha ánh xạ liên kết dùng đường đi ngắn nhất hoặc multi-path để hoàn thiện. Đồ án này đi theo hướng đó, với cải tiến là sử dụng học tăng cường để học chính sách xếp hạng và chọn ứng viên thay cho heuristic cứng.

\subsection{Hàm mục tiêu và độ đo hiệu năng}

Đối với mỗi VNR được nhúng thành công, \textbf{doanh thu} (revenue) và \textbf{chi phí} (cost) được định nghĩa:
\begin{equation}
    \mathcal{R}(\text{VNR}) = \sum_{n^v \in N^v} c(n^v) + \sum_{e^v \in E^v} b(e^v),
    \label{eq:revenue}
\end{equation}
\begin{equation}
    \mathcal{C}(\text{VNR}) = \sum_{n^v \in N^v} c(n^v) + \sum_{e^v \in E^v} b(e^v) \cdot \ell(f_E(e^v)),
    \label{eq:cost}
\end{equation}
trong đó $\ell(\cdot)$ là tổng số chặng vật lý của tập đường đi ánh xạ liên kết ảo. Hai độ đo hiệu năng phổ biến nhất là:
\begin{itemize}
\item \textbf{Tỉ lệ chấp nhận} (Acceptance Ratio - AR): tỉ số giữa số VNR được nhúng thành công và tổng số VNR đến trong khoảng thời gian khảo sát.
\item \textbf{Tỉ số doanh thu trên chi phí} (Revenue-to-Cost - R/C): được định nghĩa cho mỗi VNR là $\mathcal{R}(\text{VNR}) / \mathcal{C}(\text{VNR})$, đo mức độ ``tiết kiệm tài nguyên'' khi nhúng. Ở mức hệ thống, R/C tổng được tính bằng tỉ số của tổng doanh thu và tổng chi phí trên toàn bộ VNR đã nhúng.
\end{itemize}

Ngoài AR và R/C, các độ đo bổ sung gồm thời gian chạy trung bình mỗi VNR, độ trễ trung bình của các đường đi, và độ ổn định của reward trong quá trình huấn luyện. Đồ án này sử dụng đồng thời cả bốn độ đo để so sánh.

\subsection{Phát biểu bài toán dưới dạng ILP}

Để có hình thức hoá đầy đủ, đặt biến nhị phân $x_{n^v, n^s} \in \{0,1\}$ bằng $1$ nếu nút ảo $n^v$ được ánh xạ tới nút vật lý $n^s$, và biến $y^{u^v v^v}_{u^s v^s} \in [0, 1]$ biểu diễn tỉ lệ băng thông của liên kết ảo $(u^v, v^v)$ đi qua liên kết vật lý $(u^s, v^s)$. Với mục tiêu cực đại hoá tỉ số R/C, bài toán có thể viết:
\begin{equation}
    \max \quad \frac{\mathcal{R}(\text{VNR})}{\mathcal{C}(\text{VNR})}
\end{equation}
với các ràng buộc:
\begin{align}
    \sum_{n^s \in N^s} x_{n^v, n^s} &= 1, \quad \forall n^v \in N^v, \label{eq:c1}\\
    \sum_{n^v \in N^v} x_{n^v, n^s} \cdot c(n^v) &\leq C^{\text{avail}}(n^s), \quad \forall n^s \in N^s, \label{eq:c2}\\
    \sum_{n^v \in N^v} x_{n^v, n^s} &\leq 1, \quad \forall n^s \in N^s, \label{eq:c3}\\
    \sum_{(u^s, v^s) \in E^s} y^{u^v v^v}_{u^s v^s} \cdot b(u^v, v^v) &\leq B^{\text{avail}}(u^s, v^s), \quad \forall (u^s, v^s), \label{eq:c4}\\
    \sum_{v^s : (u^s, v^s) \in E^s} (y^{u^v v^v}_{u^s v^s} - y^{u^v v^v}_{v^s u^s}) &= x_{u^v, u^s} - x_{v^v, u^s}, \quad \forall u^s, (u^v, v^v). \label{eq:c5}
\end{align}
Ràng buộc \eqref{eq:c1} đảm bảo mỗi nút ảo được ánh xạ tới đúng một nút vật lý; \eqref{eq:c2} đảm bảo không vượt CPU; \eqref{eq:c3} đảm bảo tính đơn ánh; \eqref{eq:c4} đảm bảo không vượt băng thông; và \eqref{eq:c5} là ràng buộc bảo toàn luồng (multi-commodity flow) cho ánh xạ liên kết. Như đã phân tích ở Mục~\ref{sec:dvd} của Chương~\ref{ch:de_xuat}, dạng ILP này có $|N^v| \cdot |N^s|$ biến nhị phân nút và $|E^v| \cdot |E^s|$ biến liên kết, làm cho việc giải tối ưu trở nên không khả thi với mạng vật lý cỡ trăm nút trở lên.

\section{Các hướng tiếp cận hiện tại}
\label{sec:approaches}

\subsection{Hướng heuristic}

Các thuật toán heuristic sử dụng các luật xếp hạng đơn giản dựa trên đặc trưng tô-pô như bậc của nút, NodeRank\cite{cheng2011vne_nr}, hoặc chỉ số kết hợp giữa tài nguyên còn lại và tổng băng thông của các cạnh kề. ViNE\cite{chowdhury2009vine} giải LP-relaxation của bài toán nhị phân để gợi ý ánh xạ nút, rồi làm tròn. MP-VNE\cite{mp_vne} ước lượng chi phí của từng cặp $(n^v, n^s)$ qua công thức tổng hợp giữa CPU-cost và băng thông trung bình của các láng giềng, rồi chọn các nút có chi phí ước lượng thấp nhất làm ứng viên. Pha ánh xạ liên kết phổ biến nhất là đường đi ngắn nhất theo Dijkstra hoặc Floyd, có thể mở rộng thành \textit{multi-path} cho phép chia một liên kết ảo thành nhiều đường đi vật lý.

Ưu điểm của heuristic là đơn giản, có thể chạy trên hạ tầng lớn, và thường cho kết quả chấp nhận được trong các tình huống điển hình. Nhược điểm là phụ thuộc nặng vào tri thức chuyên gia, không thích nghi với phân phối yêu cầu thay đổi, và dễ rơi vào tối ưu cục bộ.

\subsection{Hướng metaheuristic}

Các phương pháp metaheuristic gồm thuật toán di truyền (Genetic Algorithm), tối ưu hoá bầy hạt (Particle Swarm Optimization - PSO), và mô phỏng tôi luyện (Simulated Annealing). Trong bài toán VNE, mỗi cá thể/cá nhân biểu diễn một phương án ánh xạ nút hoàn chỉnh, hàm thích nghi (fitness) là chi phí nhúng. Cá thể tốt được chọn lọc qua các thế hệ, kết hợp với phép đột biến để thoát khỏi tối ưu cục bộ. MP-VNE\cite{mp_vne} sử dụng PSO với 10 hạt, 50 vòng lặp, và xác suất đột biến 10\% để cân bằng giữa khám phá và khai thác. Hạn chế chung của nhóm metaheuristic là chi phí tính toán cao (phải đánh giá nhiều phương án trong mỗi vòng lặp) và không có đảm bảo lý thuyết về chất lượng lời giải.

\subsection{Hướng học tăng cường}

Học tăng cường (Reinforcement Learning - RL) coi VNE như một quá trình quyết định tuần tự: tại mỗi bước, tác nhân quan sát trạng thái mạng vật lý và VNR hiện tại, chọn một hành động (ánh xạ một nút ảo) rồi nhận tín hiệu thưởng phản ánh chất lượng của lời giải cuối cùng. A3C-GCN\cite{yao2018a3c_gcn} là một trong những công trình tiên phong sử dụng mạng nơ-ron đồ thị để mã hoá tô-pô mạng vật lý, kết hợp với thuật toán Actor-Critic bất đồng bộ. Gần đây, FlagVNE\cite{flagvne2024} đề xuất khung làm việc tổng quát cho RL trong cấp phát tài nguyên mạng, tách bạch giữa môi trường, biểu diễn, và thuật toán huấn luyện, đồng thời chỉ ra hiệu năng vượt trội so với heuristic và phương pháp metaheuristic trên nhiều bộ benchmark.

\subsection{So sánh và động lực của đồ án}

Bảng~\ref{tab:approach-compare} tổng hợp ưu nhược của ba nhóm tiếp cận.

\begin{table}[h]
\centering
\caption{So sánh ba hướng tiếp cận VNE}
\label{tab:approach-compare}
\resizebox{\textwidth}{!}{%
\begin{tabular}{l l l l}
\hline
\textbf{Tiêu chí} & \textbf{Heuristic} & \textbf{Metaheuristic} & \textbf{Học tăng cường} \\
\hline
Chất lượng lời giải & Trung bình & Khá & Tốt (khi huấn luyện ổn định) \\
Thời gian suy diễn & Rất nhanh & Chậm & Nhanh \\
Khả năng tổng quát hoá & Thấp & Trung bình & Cao \\
Phụ thuộc tri thức chuyên gia & Cao & Trung bình & Thấp \\
Chi phí huấn luyện & Không có & Không có & Cao \\
\hline
\end{tabular}%
}
\end{table}

Phân tích trên cho thấy hướng học tăng cường có tiềm năng lớn nhất nhưng còn hai trở ngại: (i) cần huấn luyện trên lượng mẫu lớn để hội tụ, và (ii) khó kết hợp tri thức tô-pô. Đồ án đề xuất kết hợp ưu điểm của heuristic (nguồn tri thức tô-pô) và RL (khả năng học và tổng quát hoá) thông qua quy trình huấn luyện hai giai đoạn được trình bày ở Chương~\ref{ch:de_xuat}.

\section{Học tăng cường sâu}
\label{sec:rl}

\subsection{Quá trình quyết định Markov}

Học tăng cường được hình thức hoá qua \textit{Quá trình quyết định Markov} (Markov Decision Process - MDP), một bộ năm
\[
    \mathcal{M} = (\mathcal{S}, \mathcal{A}, \mathcal{P}, \mathcal{R}, \gamma),
\]
trong đó $\mathcal{S}$ là không gian trạng thái, $\mathcal{A}$ là không gian hành động, $\mathcal{P}(s' \mid s, a)$ là xác suất chuyển trạng thái, $\mathcal{R}(s, a)$ là hàm phần thưởng, và $\gamma \in [0, 1]$ là hệ số chiết khấu. Một \textit{chính sách} $\pi(a \mid s)$ là phân phối xác suất trên hành động cho mỗi trạng thái. Mục tiêu của RL là tìm chính sách $\pi^*$ tối đa hoá \textit{kỳ vọng lợi ích tích luỹ}\cite{sutton2018rl}:
\begin{equation}
    J(\pi) = \mathbb{E}_{\tau \sim \pi}\left[ \sum_{t=0}^{T} \gamma^t \mathcal{R}(s_t, a_t) \right],
\end{equation}
trong đó $\tau = (s_0, a_0, s_1, a_1, \ldots)$ là quỹ đạo tương tác.

Trong bối cảnh VNE, một episode tương ứng với quá trình nhúng một VNR. Trạng thái mô tả mạng vật lý hiện tại và VNR đang xét, hành động là ánh xạ một nút ảo tới một nút vật lý hoặc xếp hạng các nút/liên kết ảo, và phần thưởng được cấp khi VNR đã được ánh xạ hoàn chỉnh và đo bằng tỉ số R/C.

\subsection{Policy gradient và Actor-Critic}

\textit{Policy gradient} là họ thuật toán RL học chính sách trực tiếp bằng cách tham số hoá $\pi_\theta(a \mid s)$ với tham số $\theta$ và cập nhật theo gradient của $J(\theta)$. Định lý policy gradient\cite{sutton2018rl} chỉ ra:
\begin{equation}
    \nabla_\theta J(\theta) = \mathbb{E}_{\tau \sim \pi_\theta}\left[ \sum_t \nabla_\theta \log \pi_\theta(a_t \mid s_t) \cdot G_t \right],
\end{equation}
với $G_t = \sum_{k=t}^{T} \gamma^{k-t} \mathcal{R}(s_k, a_k)$. REINFORCE\cite{williams1992reinforce} là phiên bản đơn giản nhất, ước lượng kỳ vọng bằng trung bình mẫu Monte Carlo. Hạn chế chính của REINFORCE là phương sai gradient cao. Để giảm phương sai mà không gây bias, có thể trừ một \textit{baseline} $b(s)$ độc lập với hành động:
\begin{equation}
    \nabla_\theta J(\theta) = \mathbb{E}\left[ \sum_t \nabla_\theta \log \pi_\theta(a_t \mid s_t) \cdot (G_t - b(s_t)) \right].
\end{equation}
Phương án phổ biến nhất là dùng hàm giá trị $V_\phi(s)$ làm baseline, dẫn đến kiến trúc \textit{Actor-Critic} với hai mạng song song: \textit{actor} là chính sách $\pi_\theta(a\mid s)$ quyết định hành động, còn \textit{critic} là hàm giá trị $V_\phi(s)$ ước lượng chất lượng trạng thái và được học bằng cách hồi quy về lợi ích thực tế (mất mát bình phương). Lợi thế $A(s,a)=G_t - V_\phi(s)$ thay cho $G_t$ giúp giảm phương sai gradient. Đồ án này sử dụng baseline phê bình $V_\phi(s)$ học được (kiến trúc actor-critic) thay cho baseline trung bình đơn giản, nhằm giảm phương sai mạnh hơn trong các kịch bản phần thưởng cấp ở cuối episode.

\subsection{Tối ưu hoá chính sách gần kề (PPO)}

\textit{Proximal Policy Optimization} (PPO)\cite{schulman2017ppo} là một thuật toán actor-critic cải tiến, cho phép tái sử dụng dữ liệu rollout qua nhiều vòng gradient mà vẫn ổn định. PPO thay mục tiêu gradient chính sách bằng một \textit{hàm mục tiêu thay thế có cắt} (clipped surrogate) dựa trên tỉ số tầm quan trọng $\rho_t(\theta) = \pi_\theta(a_t\mid s_t)/\pi_{\theta_{\text{old}}}(a_t\mid s_t)$:
\begin{equation}
    L^{\text{CLIP}}(\theta) = \mathbb{E}_t\Big[ \min\big( \rho_t(\theta)\,\hat{A}_t,\; \mathrm{clip}(\rho_t(\theta), 1-\epsilon, 1+\epsilon)\,\hat{A}_t \big) \Big],
    \label{eq:ppo-clip}
\end{equation}
với $\epsilon$ là biên cắt (thường $0{,}2$). Phép cắt ngăn chính sách trôi quá xa khỏi chính sách cũ trong một bước cập nhật, đóng vai trò một ràng buộc vùng tin cậy mềm. Đồ án sử dụng kiến trúc actor-critic làm phương pháp tinh chỉnh chính, đồng thời cài đặt biến thể PPO có cắt để so sánh (Chương~\ref{ch:thuc_nghiem}).

\subsection{Plackett-Luce sampling}

Trong bài toán VNE, tác nhân cần sinh ra một \textit{hoán vị} (thứ tự xử lý nút/liên kết ảo) thay vì chỉ một hành động đơn lẻ. \textit{Mô hình Plackett-Luce}\cite{plackett1975ranking} biểu diễn phân phối xác suất trên hoán vị thông qua điểm số $s_1, \ldots, s_n$:
\begin{equation}
    P(\pi \mid \mathbf{s}) = \prod_{k=1}^{n} \frac{\exp(s_{\pi(k)})}{\sum_{j \in \mathcal{R}_k} \exp(s_j)},
    \label{eq:pl}
\end{equation}
với $\mathcal{R}_k$ là tập các phần tử chưa được chọn ở các bước trước. Sampling thực hiện bằng cách softmax-sample phần tử đầu tiên, loại khỏi tập, rồi lặp lại. Log-likelihood của hoán vị có thể tính được trực tiếp, cho phép kết hợp với REINFORCE để cập nhật gradient.

\section{Mạng nơ-ron đồ thị}
\label{sec:gnn}

\subsection{Cơ chế truyền tin}

Mạng nơ-ron đồ thị (Graph Neural Network - GNN) học biểu diễn của các nút bằng cơ chế \textit{truyền tin} (message passing). Tại mỗi lớp $\ell$, vector biểu diễn $h_v^{(\ell)}$ của mỗi nút $v$ được cập nhật từ biểu diễn của các láng giềng:
\begin{equation}
    h_v^{(\ell+1)} = \text{UPDATE}^{(\ell)}\left( h_v^{(\ell)},\; \text{AGGREGATE}^{(\ell)}\left( \{h_u^{(\ell)} : u \in \mathcal{N}(v)\} \right) \right).
\end{equation}
Sau $L$ lớp truyền tin, mỗi nút ``thấy'' được thông tin từ các láng giềng trong bán kính $L$ chặng. Các biến thể GNN khác nhau chủ yếu ở hai hàm AGGREGATE và UPDATE.

\subsection{Graph Convolutional Network}

\textit{Graph Convolutional Network} (GCN)\cite{kipf2017gcn} sử dụng phép trung bình có trọng số trên láng giềng:
\begin{equation}
    H^{(\ell+1)} = \sigma\left( \tilde{D}^{-1/2} \tilde{A} \tilde{D}^{-1/2} H^{(\ell)} W^{(\ell)} \right),
    \label{eq:gcn}
\end{equation}
trong đó $\tilde{A} = A + I$ là ma trận kề có self-loop, $\tilde{D}$ là ma trận đường chéo độ bậc, $H^{(\ell)} \in \mathbb{R}^{|V| \times d_\ell}$ là ma trận biểu diễn, $W^{(\ell)}$ là tham số học được, và $\sigma$ là hàm kích hoạt phi tuyến (đồ án dùng ReLU). Đồ án sử dụng GCN hai lớp với $H^{(0)} = X$ là ma trận đặc trưng đầu vào của các nút vật lý. Ma trận kề $A$ được trọng số hoá theo tỉ lệ băng thông còn lại trên dung lượng ban đầu, cho phép mã hoá thông tin trạng thái tài nguyên động vào quá trình truyền tin.

\subsection{Graph Attention Network}

\textit{Graph Attention Network} (GAT)\cite{velickovic2018gat} thay phép trung bình đồng đều bằng phép tổ hợp có trọng số attention học được:
\begin{equation}
    h_v^{(\ell+1)} = \sigma\left( \sum_{u \in \mathcal{N}(v) \cup \{v\}} \alpha_{vu}^{(\ell)} \cdot W^{(\ell)} h_u^{(\ell)} \right),
\end{equation}
với hệ số attention
\begin{equation}
    \alpha_{vu}^{(\ell)} = \frac{\exp\left( \text{LeakyReLU}(\mathbf{a}^\top [W h_v \| W h_u]) \right)}{\sum_{u' \in \mathcal{N}(v) \cup \{v\}} \exp\left( \text{LeakyReLU}(\mathbf{a}^\top [W h_v \| W h_{u'}]) \right)}.
\end{equation}
GAT có khả năng biểu diễn cao hơn GCN nhờ trọng số học được, đặc biệt khi đồ thị không đồng nhất về vai trò các láng giềng. Tuy nhiên, GCN có ưu điểm về số tham số ít, tốc độ huấn luyện nhanh, và đủ biểu diễn cho các tô-pô mạng vật lý nhỏ - vừa trong VNE. Đồ án này lựa chọn GCN làm bộ mã hoá chính để giữ kiến trúc nhỏ gọn (toàn bộ chính sách chỉ khoảng 20{,}000 tham số), đồng thời thảo luận lại các đánh đổi với GAT trong phân tích ablation ở Chương~\ref{ch:thuc_nghiem}.

\section{Cơ chế attention cho lựa chọn ứng viên}
\label{sec:attention}

Cơ chế \textit{scaled dot-product attention}\cite{vaswani2017attention} là khối toán học chuẩn cho việc ``một truy vấn chọn từ nhiều khoá''. Cho truy vấn $q \in \mathbb{R}^d$ và tập khoá $K = (k_1, \ldots, k_n)$, điểm số cho mỗi khoá là
\begin{equation}
    \text{score}(q, k_i) = \frac{q^\top k_i}{\sqrt{d}}.
\end{equation}
Hệ số $\sqrt{d}$ giúp ổn định gradient khi $d$ lớn. Trong bài toán VNE, mỗi nút ảo đóng vai trò truy vấn, mỗi nút vật lý ứng viên đóng vai trò khoá, và điểm số sau khi softmax biểu diễn xác suất chọn từng ứng viên. Đồ án mở rộng cơ chế này bằng cách kết hợp một nhánh MLP dư (residual) để học các tương tác phi tuyến phức tạp giữa đặc trưng của nút ảo và nút vật lý, và sử dụng \textit{feasibility mask} cứng để loại bỏ các ứng viên không thoả mãn ràng buộc CPU ngay tại bước tính điểm. Chi tiết kiến trúc được trình bày ở Chương~\ref{ch:de_xuat}.

% Kết chương
Chương này đã hệ thống hoá các khái niệm và mô hình cần thiết cho phương pháp đề xuất, bao gồm mô hình hoá VNE và phát biểu ILP, khảo sát ba hướng tiếp cận hiện hành, nền tảng học tăng cường (MDP, gradient chính sách, actor-critic, PPO, Plackett-Luce), và các khối GNN và attention. Chương tiếp theo sử dụng các nền tảng này để thiết kế kiến trúc và thuật toán huấn luyện của CARL-VNE.

\end{document}
